% !tp-example Basic usage
\documentclass{article}
\begin{document}
\wcTpFillBlankQuestion[mode=closed]{
  \begin{description}
    \item[{Title}] Match the answers
    \item[{Prompt}] Choose an answer for each blank.
    \item[{Answers}] \begin{enumerate}
      \item Paris
      \item \rawhtml|<svg xmlns="http://www.w3.org/2000/svg" width="48" height="40" viewBox="0 0 48 40" role="img" aria-label="Triangle"><path d="M24 3 45 37H3Z" fill="none" stroke="currentColor" stroke-width="3"></path></svg>|
      \item \wcImg*[src=/docs/medias/logos/logo-tp.svg,alt=tp-components logo,width=40,height=40]{}
    \end{enumerate}
    \item[{Form}] The capital of France: \wcTpBlank*[name=capital,aria-label=Capital of France]{}
     A shape with three sides: \wcTpBlank*[name=shape,aria-label=Three-sided shape]{}
     The library logo: \wcTpBlank*[name=logo,aria-label=Library logo]{}
    \item[{Feedback}] Match each answer to its description.
    \item[{Solution}] Paris; the triangle; the tp-components logo.
  \end{description}
}
\end{document}
